\par\SourceLeaf{127}
It follows that liability II. too can be valued in money with a fair degree of reliability, and that, to carry out these calculations for I. and II., the year and day of birth must be obtained for
\begin{quote}
each of the widows now living,\par
their children under 20 years of age, where there are any, as in the cases of Mrs Hofr. M., Mrs G. J. R. M., and Mrs Prof. H.,\par
the members who are now married,\par
their wives.
\end{quote}
\noindent\hspace*{\parindent}Ad III. Most questionable, however, must seem the undertaking of expressing in figures the present monetary value of the fund's liability towards future members \textit{in saecula saeculorum}, under, of course, the present statutes and at the present levels of pensions and contributions. And yet it is \textit{necessary} to be put in a position to form at least an approximate idea of this, since the very question is \textit{to assess the stability}, not of a widows' fund whose arrangements may subsequently be altered according to circumstances, but \textit{of our widows' fund under its present arrangements.} Naturally, one must not forget that the calculation remains dependent on certain elements which can in part already be estimated only approximately and can in part undergo very considerable changes over time. I shall name two such elements: the level of the interest rate, and the average number of new members joining each year.

For an institution which is, and is to remain, based only to a very small extent on continuing contributions, and chiefly on income from capital, the level of the interest rate evidently bears an exact (inverse) relation to the amount of capital required. Thus, if, for example, at some point a gift of capital of 70000 Thl. just sufficed to sustain a society admitting on average the same number of new members each year at an interest rate of 4 p.c., a fall in the interest rate to \(3\frac12\) p.c. would require the capital to be increased to 80000 Thl. I regard this circumstance as precisely the least significant \textit{in relation to the difficulty of the assessment}. The assessment can do no more than make the size of its influence clear; from this it follows of itself that provision must already be made for the capital to receive a suitable average annual increase, to counter the gradual but unfailing fall in the interest rate over time.

It is equally evident that the amount of capital required will be exactly proportional to the average number of new members joining each year (\textit{ceteris paribus}). For the time being we can base ourselves only on our own experience, which extends over 100 years, and in which the more recent and most recent periods must naturally guide our judgment above all. His Magnificence rightly observes that the increase in the number of professors arising from increased scholarly needs and demands must have a substantial influence on the continued existence of professors' widows' funds of this kind, which are chiefly based on capital. It is therefore quite possible that the Göttingen results, for example over the last 30 or 40 years, cannot provide an entirely correct measure for the future, especially for the future of later centuries. But this uncertainty lies in the nature of the changeability of all human affairs; its consequences emerge \textit{gradually}, and can be met only by vigilance that never falls asleep. In our case, therefore, one calculates the capital required to meet liability III. on a lasting basis according to the best of our present knowledge, does not forget that this very uncertainty requires a somewhat larger capital to be available, \textit{repeats} the calculation continually at definite, not overly long
\SourceLeaf{128}
intervals, for example every 5 or 10 years, always making use of the newly acquired experience, and regards the surplus as partly available for disposal only when it has actually \textit{increased}.

But even apart from these two circumstances, or, in other words, even if one takes as a basis a definite interest rate and a definite average number of new members joining each year, the difficulty of the estimate still seems almost insuperable, since the ages of the spouses can stand in the most varied relations, not even to mention the remarriage of those who have been widowed. In this respect, therefore, the assessor seems to stand on roughly the same footing as someone who would have had to examine in advance the plan of one of the many widows' funds which, established without rigorous regard to the ages of the married couples joining, have almost all come to ruin. (If University Councillor K. believes that there will have been no lack of \textit{calculation} in all these funds either, he is undoubtedly right; but if he wishes to infer from this that the \textit{calculus of probabilities} has no sound foundation, he is wrong. There are indeed many words to which different people attach different meanings, and likewise words with a very definite scientific meaning under which everyday usage often lumps together very disparate things. Such is the case with the expression calculus of probabilities. Understood in its strict sense, there can be no question of applying it in any case where the necessary foundations are lacking. In establishing the arrangements of all the failed widows' funds, \textit{no use at all} was made of the rigorous calculus of probabilities, but only of vague impressions. I state this here merely as a fact, not as a reproach, since a basis in the calculus of probabilities was in fact impossible for the very reason that all the necessary conditions for it were lacking.) In the present case, however, although it is impossible to determine a final result by the calculus of probabilities from the individual elements, precisely because these elements are lacking, the fund's own hundred years of experience, if used in the right way, offer a rich store on which to base it. This experience must therefore first be collected and put in order. I presuppose the establishment of a book in which all members without exception, from the society's first foundation onwards, are entered in chronological order of joining, together with all information relevant to the purpose in question. This experience would certainly provide still much more fruitful material if the year and day of birth had also been recorded for all the persons concerned, namely for the professor joining, for his wife if he is already married, or if and whenever he marries after joining, and finally for the minor children living at the member's death. But all these things are lacking, and could now in most cases be supplied only for the member himself; little or no benefit could be derived from these isolated particulars. Nevertheless, what can undoubtedly still be brought together is of the greatest value. I mean, for each individual member,

1. The \textit{terminus a quo} and \textit{ad quem} of the contributions paid.

2. The \textit{terminus a quo} and \textit{ad quem} of the widow's pension received, in the cases where such a pension occurred.

3. The \textit{terminus a quo} and \textit{ad quem} of the orphan's pension received, where minor children remained after the death of the father or mother.

But the \textit{amounts} of money contributed by the members and drawn by the widows and orphans need not be extracted as well.

This, then, would be the \textit{third requirement}, which must indispensably be obtained before any calculations begin. I am entirely unfamiliar with the arrangement of the widows' fund's archive, and therefore do not know whether specially established books may already exist from which this material can be obtained with certainty, completeness
\SourceLeaf{129}
and ease. In any case, however, the 105 annual accounts, which have undoubtedly been preserved, could serve this purpose. Only after taking a closer look at the available papers myself could I say whether, and to what extent, I can promise my own help in this extraction.

For each individual member to be entered in this book, the data under the respective headings allow one to calculate how much he has paid to the society and how much it has paid to his surviving dependants at the present rates, and what the monetary value of this would have been, reckoned back at a definite interest rate to the time when he joined. Naturally, this varies greatly between individuals, being positive for some and negative for others. But by the eternal laws, the \textit{mean} of a large number is an element which can again be taken as a mean for the future, where there is no reason to suppose essential changes in the general circumstances. The table itself will already be able to give instructive indications on this point if one groups the whole and compares, for example, the mean for the first half with that for the second, or the first, second, and third thirds with one another. It will thus emerge how large is the monetary value of the liability which accrues to the society \textit{on average} when a new member joins. If, in accordance with what has already been remarked above, one has also made a plausible assumption about the average number joining annually, one can calculate, for a definite interest rate, how large the capital must be whose interest income can cover this perpetually continuing liability III. The three capitals for I., II., and III., added together and compared with the society's actual assets, will then show, as precisely as the nature of the subject permits, whether under its present arrangements its stability is \textit{more} than assured, or only just adequate, or whether its instability follows, so that its ruin must decidedly be expected sooner or later. Accordingly, the appropriate measures would then have to be considered and initiated in the different cases.

These are the main outlines of the plan according to which, in my opinion, a thorough examination and the drawing up of a balance sheet could and should be carried out. It is a substantial piece of work, which I shall, however, gladly undertake if desired. The shortness of the time will excuse my presenting this sketch only in a rough and hastily written form. If such a work is carried out now, it remains in any case urgently desirable, as already remarked above, that a new balance sheet should be drawn up repeatedly at definite intervals in future. This would then be much easier than the first time if a book such as I mentioned above were established, at least from now on, with all the \textit{date headings}, and regularly completed and continued.

I shall now also add my opinion on a couple of other points touched on in the other votes.

I am not in favour of repealing the provision of the statutes which excludes unsalaried professors. In relation to the widows' fund, the University as a corporation would always have to wish urgently that cases in which the title of professor is conferred on someone employed in a school or church should remain very rare. Apart from such a case, however, someone with neither a salary nor assets is unlikely to be so imprudent as to marry, and thus the entire provision will seldom, perhaps never, have any effect. But it is possible that an unsalaried professor who, feeling himself near death, had a bride whom he would certainly not otherwise have married before obtaining a salary might, if entry into the widows' fund were open to him, be tempted to burden the fund by a hasty marriage ceremony.

As long as people act \textit{bona fide}, unsalaried professors must rather regard that clause as
\SourceLeaf{130}
a fair provision to their advantage, which spares them the alternative of either contributing to the widows' fund during the very time when they receive nothing from the University, or later, once they receive a salary, having to make double back payments for the entire period of their unsalaried professorship.

\Needspace{32\baselineskip}
My second remark concerns the interest rate, on which I cannot entirely agree with what is said in the P. M. of U. R. O. Rather, the calculation by which the present interest rate is determined as \(4\tfrac1{10}\) per cent seems to me in part illusory. I shall explain myself by an example. According to today's price list, Austrian \(4\tfrac12\) per cent securities stand at \(103\tfrac58\). When buying from a banker, one will certainly have to pay \textit{more than} 104 in all; but I shall keep to 104 alone. Thus, on the money paid in, one actually receives only \(4\tfrac{17}{32}\), or not quite \(4\tfrac13\) per cent interest. It therefore takes at least 12 years before one can even say that one has actually enjoyed 4 per cent interest. Now, however, \textit{very large} sums of these securities are drawn each year and repaid at par. If the drawing occurs after only 2 years, one has actually enjoyed only \(4\tfrac23\) per cent in total, or \(2\tfrac13\) per cent for one year, leaving aside the costs attached to each collection. It is also always a dangerous circumstance for the holder of such a security that, if the drawing affecting him does not come to his attention and he therefore misses collecting at the proper time, he may suffer a very substantial loss. For the widows' fund, the banker from whom the securities were bought will probably always exercise the necessary vigilance, since every transaction that arises brings him a profit; but he will hardly assume actual responsibility for every loss caused by a possible oversight. In this connection, I shall therefore not fail to report here that, in the drawing of \textit{one and a half million gulden} of these securities which took place in Vienna eight days ago today, one of the widows' fund's bonds was also drawn, namely the one listed on page IX. of the accounts under No. 52, Litr. P No. 15472. I owe my ability to make this report only to the chance circumstance that today, while these very accounts are in my hands, I found a notice of the drawing in a newspaper and therefore noted down the designation of the number drawn, so that I could compare it at home with the widows' fund's list of capital holdings. With this report I shall bring this long exposition to a close.

\noindent 9. January 1845.\hfill\textsc{Gauss.}

